\documentclass[11pt,a4paper]{article}
\usepackage[margin=2.55cm]{geometry}
\usepackage[T1]{fontenc}
\usepackage{lmodern}
\usepackage{amsmath,amssymb,booktabs,graphicx,caption,microtype}
\usepackage[colorlinks=true,linkcolor=black,urlcolor=black]{hyperref}
\usepackage{fancyhdr}
\setlength{\parindent}{0pt}
\setlength{\parskip}{4pt}
\captionsetup{font=small,labelfont=bf}
\pagestyle{fancy}\fancyhf{}\fancyhead[L]{\footnotesize From Yield Curves to Bond Prices}\fancyhead[R]{\footnotesize November 2026}\fancyfoot[C]{\footnotesize\thepage}\renewcommand{\headrulewidth}{0.3pt}
\newcommand{\figdir}{figures/}
\begin{document}
\begin{center}
{\LARGE\bfseries From Yield Curves to Bond Prices\par}
\vspace{3pt}
{\large A small numerical study with the Vasicek model\par}
\vspace{12pt}
Kelvin Schouten\par
November 2026
\end{center}
\vspace{3pt}
\begin{abstract}
This note explores how a simple interest-rate model can be fitted to a yield curve and used for pricing. I use Dutch central bank (DNB) spot yields for 31 October 2022, calibrate the Vasicek model, simulate short rates, and value a European option on a zero-coupon bond. I also fit a Nelson--Siegel curve to the same observations. Nelson--Siegel gives a slightly smaller in-sample fitting error, while Vasicek provides a stochastic model that can be used for risk-neutral pricing. The aim is to understand what the models do and where their limitations become visible, rather than to build a forecasting model.
\end{abstract}

\section{Starting with a yield curve}
An observed yield curve tells us how the interest rate varies with maturity. I start with the DNB dataset for 31 October 2022, containing continuously compounded spot yields $R(T)$ at integer maturities $T=1,\ldots,100$ years. Yields are reported in percent and converted to decimals in the calculations.

From the spot rates I calculate zero-coupon discount factors and one-year forward rates:
\begin{equation}
P(0,T)=e^{-TR(T)},\qquad F(T,T+1)=\log\!\left(\frac{P(0,T)}{P(0,T+1)}\right).
\end{equation}
The spot curve rises initially, reaches $3.085\%$ at 12 years, and then falls to $1.558\%$ at 100 years. The shape suggests a simple question: can a model driven by only one stochastic interest rate reproduce this pattern?
\begin{figure}[htbp]\centering
\includegraphics[width=.74\linewidth]{\figdir VRM_P2_1_Spot.png}
\caption{Observed DNB spot yields on 31 October 2022.}\label{fig:spot}
\end{figure}

\section{Fitting a stochastic short-rate model}
The Vasicek model assumes that the short rate follows an Ornstein--Uhlenbeck process under a pricing measure $Q$:
\begin{equation}
dr_t=a(b-r_t)\,dt+\sigma\,dW_t^Q,\qquad a>0.
\end{equation}
Here $a$ determines how quickly the rate moves towards $b$, while $\sigma$ controls the size of random shocks. This gives a convenient closed-form expression for the model-implied spot yield:
\begin{align}
R^{\mathrm{Vas}}(T)&=R_\infty+\frac{1-e^{-aT}}{aT}(r_0-R_\infty)
+\frac{\sigma^2(1-e^{-aT})^2}{4a^3T},\\
R_\infty&=b-\frac{\sigma^2}{2a^2}.
\end{align}
I estimate $(a,b,\sigma,r_0)$ by minimizing the equally weighted squared differences between the 100 observed yields and the model yields. The initial short rate $r_0$ is fitted as well, since the dataset starts at one year rather than at maturity zero.

\begin{table}[htbp]\centering
\caption{Vasicek calibration to the DNB yield curve.}\label{tab:vas}
\begin{tabular}{ccccc}\toprule
$a$ & $b$ & $\sigma$ & $r_0$ & Yield RMSE\\\midrule
0.10484 & 0.07099 & 0.03628 & 0.02359 & 5.26 bp\\\bottomrule
\end{tabular}
\end{table}
The model picks up the hump in the observed curve, but does not fit every maturity exactly. Its maximum absolute yield error is 13.81 basis points. One limitation is that a single factor constrains how much flexibility the curve can have.
\begin{figure}[htbp]\centering
\includegraphics[width=.73\linewidth]{\figdir VRM_P2_2_Vasicek_Fit.png}
\caption{Observed yields and the fitted Vasicek curve.}\label{fig:vas}
\end{figure}

\section{What happens when rates evolve?}
Once calibrated, the model can be simulated. Rather than using an Euler approximation, I use the exact monthly transition:
\begin{equation}
r_{t+\Delta}=b+(r_t-b)e^{-a\Delta}
+\sigma\sqrt{\frac{1-e^{-2a\Delta}}{2a}}Z,\qquad Z\sim N(0,1).
\end{equation}
Figure~\ref{fig:paths} shows six example paths over five years using $\Delta=1/12$. Although each path is different, the conditional expected rate moves towards the fitted level $b=7.10\%$. After five years the model implies $E^Q[r_5]\approx4.29\%$ and $\operatorname{SD}^Q(r_5)\approx6.39\%$.

\begin{figure}[htbp]\centering
\includegraphics[width=.73\linewidth]{\figdir VRM_P2_3_Vasicek_Paths.png}
\caption{Six simulated monthly short-rate paths and the conditional mean under $Q$.}\label{fig:paths}
\end{figure}

These are \emph{risk-neutral} paths, not forecasts of future observed rates. The parameters were chosen to match current prices under $Q$, not to estimate historical dynamics under the real-world measure $P$. Vasicek also allows negative rates, which is mathematically convenient but may be undesirable in some applications.

\section{A simple bond-option experiment}
I next consider European call and put options expiring in $T_1=0.75$ years on a zero-coupon bond maturing in $T_2=1$ year. Both options have strike $K=0.98$. Their time-zero values satisfy
\begin{align}
C_0&=E^Q\!\left[e^{-\int_0^{T_1}r_s ds}\big(P(T_1,T_2)-K\big)^+\right],\\
U_0&=E^Q\!\left[e^{-\int_0^{T_1}r_s ds}\big(K-P(T_1,T_2)\big)^+\right].
\end{align}
In Vasicek, the short rate at option expiry and its accumulated integral are jointly Gaussian. I draw 200,000 pairs using their known mean and covariance, evaluate the bond price at option expiry, and discount each payoff. A fixed random seed makes this calculation repeatable.
\begin{table}[htbp]\centering
\caption{Monte Carlo option prices per unit bond face value.}\label{tab:options}
\begin{tabular}{lcc}\toprule
Option & Estimated price & Monte Carlo standard error\\\midrule
Call & 0.013048 & 0.0000161\\
Put & 0.000104 & 0.00000159\\\bottomrule
\end{tabular}
\end{table}
As a numerical check, put--call parity requires
\begin{equation}
C_0-U_0=P(0,T_2)-KP(0,T_1).
\end{equation}
The simulated left-hand side is 0.012944, compared with 0.012955 from the bond-pricing formula. The difference is small relative to the simulation error. This does not prove that the model describes market prices correctly, but it helps verify that the numerical pricing procedure is internally consistent.

\section{Would a different curve fit better?}
For comparison I fit the four-parameter Nelson--Siegel curve to the same 100 observations:
\begin{equation}
R^{\mathrm{NS}}(T)=\beta_0+\beta_1\frac{1-e^{-\lambda T}}{\lambda T}
+\beta_2\left(\frac{1-e^{-\lambda T}}{\lambda T}-e^{-\lambda T}\right).
\end{equation}
Its fitted parameters are $\beta_0=0.011274$, $\beta_1=0.012541$, $\beta_2=0.044339$, and $\lambda=0.144762$.
\begin{figure}[htbp]\centering
\includegraphics[width=.73\linewidth]{\figdir VRM_P2_5_Nelson_Siegel_Comparison.png}
\caption{Observed yields compared with Nelson--Siegel and Vasicek.}\label{fig:compare}
\end{figure}
\begin{table}[htbp]\centering
\caption{In-sample yield-curve fitting errors.}\label{tab:compare}
\begin{tabular}{lcc}\toprule
Model & RMSE (bp) & Maximum absolute error (bp)\\\midrule
Vasicek & 5.26 & 13.81\\
Nelson--Siegel & 4.95 & 13.13\\\bottomrule
\end{tabular}
\end{table}
Nelson--Siegel has a slightly better in-sample fit. It is useful for describing the curve's level, slope, and curvature, but the static formula alone does not specify the risk-neutral rate dynamics needed for derivative pricing. The comparison therefore depends on what we want to do: fit today's yields or describe a model for future rate-dependent payoffs.

\section{Discussion}
This small experiment connects three tasks that initially seem separate: observing a yield curve, choosing a stochastic model, and pricing an option. The calibrated Vasicek model captures the main shape of the DNB curve and makes simulation and bond-option valuation straightforward. Nelson--Siegel fits the same cross section slightly better, but does not by itself replace a stochastic pricing model.

There are important limits. Both fits use the same date and maturities on which they are evaluated; this is not an out-of-sample test. The 100-year range includes long maturities, where a parametric curve can be especially restrictive. Furthermore, calibrating under $Q$ does not establish how interest rates will behave under $P$. A natural next step would be to study parameter sensitivity or compare the fit across dates, rather than concluding that one model is universally better.

\end{document}
